\documentclass[../Main.tex]{subfiles}

\begin{document}
\section{Inner Products}
\begin{definition}{Inner product}
    For a vector space $V$, an \underline{inner product} is a positive-definite symmetric bilinear form (for $\F = \R$) or a Hermitian sesquilinear form (if $\F = \C$) $\varphi : V \times V \mapsto \F$. The pair $(V, \varphi)$ is an \underline{inner product space}.
\end{definition}
We denote $\varphi(u, v) = \inn{u}{v}$.

\begin{lemma}[Cauchy-Schwarz Inequality]
    For vectors $u, v$ in an inner product space,
    \begin{equation*}
        \abs{\inn{u}{v}} = \norm{u} \norm{v}
    \end{equation*}
    \label{lemCauchySchwarz}
\end{lemma}
\begin{proof}
    Consider the discriminant of $\norm{tu + v}^2$ for $t \in \F$.
\end{proof}
\begin{corollary}[Triangle Inequality]
    For all $u, v$ in an inner product space,
    \begin{equation*}
        \norm{u + v} \leq \norm{u} + \norm{v}
    \end{equation*}
    \label{corTriangleIneq}
\end{corollary}
\begin{proof}
    Evaluate $\norm{u + v}^2$
\end{proof}
\section{Orthogonality}
\begin{definition}{Orthogonal set}
    For an inner product space $(V, \varphi)$, a set $E \subseteq V$ is \underline{orthogonal} if, for distinct $d, e \in E$, $\inn{d}{e} = 0$. If also $\norm{e} = 1$ for all $e\in E$ then the set is \underline{orthonormal}.
\end{definition}
% The rest of this section is quite boring.
\section{Maps of Inner Product Spaces}
\subsection{Adjoint Map}
\begin{theorem}[Adjoint map]
    Let $V, W$ be finite-dimensional inner product spaces, $\alpha \in \L(V, W)$. There exists a unique linear map $\alpha^* : W \mapsto V$ such that:
    \begin{equation*}
        \inn{\alpha(v)}{w}_W = \inn{v}{\alpha^*(w)}_V
    \end{equation*}
    and further, for orthonormal bases $B, C$:
    \begin{equation*}
        [\alpha^*]^C_B = \left([\alpha]^B_C\right)^\dagger
    \end{equation*}
    \label{thmAdjoint}
\end{theorem}
\begin{proof}
    First, define $\alpha^*$ by the transpose of the matrix of $\alpha$. Show that this satisfies the inner product property by considering vectors and matrices.

    Show uniqueness by considering another map $\beta$ and matrix $B$ satisfying the condition (the adjoint of $\alpha$). Find the elements by considering the inner product condition on the basis vectors, and show that $B = A^\dagger$ and so $\beta = \alpha^*$ as already defined.
\end{proof}
\begin{definition}{Adjoint}
    For $V, W$ finite-dimensional inner product spaces, and $\alpha \in \L(V, W)$, the \underline{adjoint} to $\alpha$ is $\alpha^*$ as defined in \thmref{thmAdjoint}.
\end{definition}
\begin{theorem}[Riesz Representation Theorem]
    For $(V, \varphi)$ a finite-dimensional real inner product space, the map:
    \begin{equation*}
        \varphi_R : V \mapsto V^*
    \end{equation*}
    where $(\varphi_R(v))(u) = \inn{u}{v}$ is an isomorphism.
    \label{thmRieszRepresentation}
\end{theorem}
\begin{proof}
    We already know that $\varphi_R$ is linear, so we just need to show the kernel is trivial. For $v \in \ker(\varphi)$, we must have $\inn{v}{v} = 0$ and so $v = \zv$.
\end{proof}
\begin{definition}{Self-adjoint}
    For a finite-dimensional inner product space $(V, \varphi)$, $\alpha \in \L(V, V)$ is \underline{self-adjoint} if $\alpha^* = \alpha$.
\end{definition}
This is equivalent to $[\alpha]^B_B$ Hermitian in any orthonormal basis.
\subsection{Isometries}
\begin{definition}{Isometry}
    For $V, W$ inner product spaces, $\alpha\in \L(V, W)$ is an \underline{isometry} if, for all $u, v \in V$,
    \begin{equation*}
        \inn{\alpha(u)}{\alpha(v)}_W = \inn{u}{v}_V
    \end{equation*}
\end{definition}
\begin{proposition}
    For $\alpha$ an invertible linear map of finite inner product spaces $\alpha : V \mapsto W$, $\alpha$ is an isometry if and only if $\alpha^{-1} = \alpha^*$.
    \label{propIsomIffInverseAdj}
\end{proposition}
\begin{proof}
    If we have the inverse condition then we can easily show $\alpha$ is an isometry.

    If instead $\alpha$ is an isometry, consider $\inn{u}{\alpha^*(\alpha(v)) - v}$ and show it is zero.
\end{proof}
This condition is equivalent to the matrix of $\alpha$ being unitary (for $\C$) or orthogonal (for $\R$).
\section{Eigenvectors and Eigenvalues}
Let $V$ be a finite-dimensional vector space over $\R$ or $\C$ and let $\alpha \in \L(V, V)$.
\begin{propositions}{
        Let $\alpha$ be self-adjoint. Then:
        \label{propsSelfAdjEVecEVal}
    }
    \item all eigenvalues of $\alpha$ are real; \label{propSelfAdjRealEVals}
    \item eigenvectors corresponding to different eigenvalues are orthogonal. \label{propSelfAdjEVecOrtho}
\end{propositions}
\begin{proof}
    First consider $\inn{\lambda v}{v}$ and move the $\lambda$ to the right by self-adjointness.

    Then consider $\lambda \inn{v}{w}$ with $v, w$ having different eigenvalues.
\end{proof}
\begin{lemma}
    For any subspace $W$ of $V$, if $\alpha(W) \leq W$ then also $\alpha^*(W^\perp) \leq W^\perp$.
    \label{lemSelfAdjSubspaces}
\end{lemma}
\begin{proof}
    Take an inner product between $w \in W$ and $\alpha^*(u)$ where $u \in W^\perp$.
\end{proof}
\begin{lemma}
    Let $\alpha$ be self-adjoint. Then $\alpha$ has a real eigenvector.
    \label{lemRealEVec}
\end{lemma}
\begin{proof}
    We must have an eigenvalue of $\alpha$ and by \thmref{propSelfAdjRealEVals} it is real. Decompose the eigenvector $v$ into real and imaginary parts, and apply $A = [\alpha]_B^B$. Take the real part to show $\Re(v)$ is an eigenvector.
\end{proof}
\begin{theorem}[Spectral Theorem for Self-Adjoint Endomorphisms]
    Let $\alpha$ be self-adjoint. Then $V$ has an orthonormal basis of real eigenvectors of $\alpha$.
    \label{thmSpectral}
\end{theorem}
\begin{proof}
    Proceed by induction.

    Use the above lemmas to do the following:
    \begin{enumerate}
        \item find a real eigenvector $v$ and eigenvalue $\lambda$;
        \item note that $\alpha$ preserves $\spn{v}^\perp$ so define $\beta$ to be the restriction to this subspace;
        \item show that $\beta$ is self-adjoint and apply the inductive hypothesis to get a basis;
        \item normalise $v$ and add it as a basis vector.
    \end{enumerate}
\end{proof}
There also exists a spectral theorem for unitary matrices. In this case, all eigenvalues are of unit modulus. There is no spectral theorem for orthogonal matrices, as illustrated by the general 2D rotation matrix.
\begin{definition}{Normal map}
    Let $V$ be a finite-dimensional complex inner product space. $\alpha \in \L(V,V)$ is \underline{normal} if $\alpha\alpha^* = \alpha^*\alpha$.
\end{definition}
Both self-adjoint and unitary maps are normal. We can in fact prove that normal maps have a spectral theorem. This is the largest class of maps with a spectral theorem.
\end{document}